% Part:sets-functions-relations
% Chapter: size-of-sets
% Section: equinumerous-sets

\documentclass[../../../include/open-logic-section]{subfiles}

\begin{document}

\olfileid{sfr}{siz}{equ}

\olsection{ગણસામ્ય}

ગણોના ``કદ''નો આપણો સહજ ખ્યાલ સાન્ત ગણો માટે સારી રીતે
કામ કરે છે. પરંતુ અનંત ગણોનું શું? કોઈ પણ કદના બે ગણોના કદની
તુલના કરવાની ઔપચારિક રીત ઘડવી હોય, તો ગણો ક્યારે સરખા કદના છે
તેની વ્યાખ્યાથી શરૂઆત કરવી યોગ્ય છે. અહીં ફ્રેગેનું ઉદાહરણ છે:
\begin{quote}
  પીરસનારને મેજ પર થાળીઓ જેટલી જ છરીઓ મૂકી છે તેની ખાતરી કરવી
  હોય, તો તેને બંનેમાંથી એકેયની ગણતરી કરવાની જરૂર નથી. તે દરેક
  થાળીની જમણી બાજુએ એક છરી એવી રીતે મૂકે કે મેજ પરની દરેક છરી
  કોઈ થાળીની જમણી બાજુએ હોય એટલું પૂરતું છે. આમ, થાળીઓ અને
  છરીઓ એકબીજા સાથે અનન્ય રીતે, અને ખરેખર એ જ અવકાશી સંબંધ
  દ્વારા, સંગત કરવામાં આવે છે. \citep[\S70]{Frege1884}
\end{quote}
આ અવતરણનો મુખ્ય વિચાર ઔપચારિક વ્યાખ્યા દ્વારા સ્પષ્ટ કરી શકાય છે:

\begin{defn}\ollabel{comparisondef}
  $A$ અને $B$ \emph{સામ્ય ગણો} છે, જે $\cardeq{A}{B}$ લખીએ,
  જો અને માત્ર જો કોઈ !!a{bijection} $f \colon A \to B$ હોય.
\end{defn}

\begin{prop}\ollabel{equinumerosityisequi}
ગણસામ્ય એ સામ્ય સંબંધ છે.
\end{prop}

\begin{proof}
આપણે બતાવવું છે કે ગણસામ્ય સ્વવાચક, સંમિત અને પરંપરિત છે.
$A, B$ અને $C$ ગણો હોય.

\emph{સ્વવાચકતા.} તદેવ નિરૂપણ $\Id{A} \colon A \to A$, જ્યાં
બધા $x \in A$ માટે $\Id{A} (x) = x$, એ !!a{bijection} છે.
તેથી $\cardeq{A}{A}$.

\emph{સંમિતતા.} ધારો કે $\cardeq{A}{B}$; એટલે કે કોઈ
!!a{bijection} $f\colon A \to B$ છે. $f$ એ !!{bijective} હોવાથી
તેનું પ્રતિવિધેય $f^{-1}$ અસ્તિત્વમાં છે અને !!{bijective} પણ છે.
તેથી $f^{-1}\colon B \to A$ એ !!a{bijection} છે અને
$\cardeq{B}{A}$.

\emph{પરંપરિતતા.} ધારો કે $\cardeq{A}{B}$ અને $\cardeq{B}{C}$;
એટલે કે !!{bijection}s $f\colon A \to B$ અને $g\colon B \to
C$ છે. ત્યારે સંયોજન $\comp{f}{g}\colon A \to C$ એ !!{bijective}
છે, તેથી $\cardeq{A}{C}$.
\end{proof}

\begin{prop}
જો $\cardeq{A}{B}$ હોય, તો $A$ એ !!{enumerable} હોય જો અને
માત્ર જો $B$ પણ હોય.
\end{prop}

\begin{editorial}
નીચેની સાબિતીમાં \olref[enm]{sec} સમાવેલો હોય, તો
\olref[enm]{defn:enumerable}ની વ્યાખ્યા, અને અન્યથા
\olref[enm-alt]{defn:enumerable}ની વ્યાખ્યા વપરાય છે.
\end{editorial}

\begin{proof}
ધારો કે $\cardeq{A}{B}$, એટલે કોઈ !!{bijection}
$f \colon A \to B$ છે, અને ધારો કે $A$ એ !!{enumerable} છે.
\oliflabeldef{sfr:siz:enm:defn:enumerable}{
  ત્યારે કાં તો $A = \emptyset$ અથવા કોઈ !!a{surjective} વિધેય
  $g\colon \PosInt \to A$ છે. જો $A = \emptyset$ હોય, તો
  $B = \emptyset$ પણ હોય (અન્યથા કોઈ !!a{element}~$y \in B$ હોત,
  પરંતુ $f(x) = y$ થાય એવો કોઈ $x \in A$ ન હોત). બીજી તરફ,
  જો $g\colon \PosInt \to A$ એ !!{surjective} હોય, તો
  $\comp{g}{f} \colon \PosInt \to B$ પણ !!{surjective} છે.
  આ જોવા $y \in B$ લો. $f$ એ !!{surjective} હોવાથી એવો
  $x \in A$ છે કે $f(x) = y$. $g$ એ !!{surjective} હોવાથી
  એવો $n \in \PosInt$ છે કે $g(n) = x$. તેથી
\[
(\comp{g}{f})(n) = f(g(n)) = f(x) = y
\]
અને આમ $\comp{g}{f}$ એ !!{surjective} છે. $\comp{g}{f}$ એ
$B$ની પરિગણના છે, તેથી $B$ એ !!{enumerable} છે.}
{
  ત્યારે કાં તો $A  = \emptyset$ અથવા કોઈ !!a{bijection}~$g$ છે
  જેનો વિસ્તાર $A$ છે અને જેનો પ્રદેશ $\Nat$ અથવા પ્રાકૃતિક
  સંખ્યાઓનો કોઈ આરંભખંડ છે. જો $A = \emptyset$ હોય, તો
  $B = \emptyset$ પણ હોય (અન્યથા કોઈ~$y \in B$ હોત, પરંતુ
  $f(x) = y$ થાય એવો કોઈ $x \in A$ ન હોત). હવે ધારો કે આપણું
  !!{bijection}~$g$ છે. ત્યારે $\comp{g}{f}$ એ વિસ્તાર~$B$ અને
  $g$ના પ્રદેશ જેવો જ પ્રદેશ (એટલે કે $\Nat$ અથવા તેનો કોઈ
  આરંભખંડ) ધરાવતું !!a{bijection} છે. તેથી $B$ એ !!{enumerable} છે.}
\sourcecorrection{OLSIZ-006}{મૂળની \texttt{equinumerous-sets.tex},
પંક્તિઓ 68--94ની બંને સશરતી શાખામાં ખાલી ગણના કિસ્સે
\(g(x)=y\) લખાયું છે, જોકે એ શાખામાં \(g\) વ્યાખ્યાયિત નથી.
અહીં પહેલેથી આપેલા સામ્ય વિધેય મુજબ \(f(x)=y\) લખ્યું છે.}

જો $B$ એ !!{enumerable} હોય, તો ઉપરની દલીલ $f$ને બદલે
!!{bijection} $f^{-1}\colon B \to A$ સાથે ફરી કરતાં $A$ એ
!!{enumerable} છે તેમ મળે છે.
\end{proof}

\begin{prob}
બતાવો કે જો $\cardeq{A}{C}$ અને $\cardeq{B}{D}$ હોય તથા
$A \cap B = C \cap D = \emptyset$ હોય, તો
$\cardeq{A \cup B}{C \cup D}$.
\end{prob}

\begin{prob}
બતાવો કે જો $A$ અનંત અને !!{enumerable} હોય, તો
$\cardeq{A}{\Nat}$.
\end{prob}

\end{document}
